\section{Related Work}
\label{related_work}

\paragraph{Deep visual object tracking.}
Modern deep trackers learn robust target representations and perform end-to-end search/matching under appearance changes, motion blur, and partial occlusion.
Recent work also targets highly dynamic objects that exhibit large deformation and non-rigid motion.
In this paper, we adopt \emph{NetTrack}~\cite{zheng2024nettrack} as a strong deep-learning baseline, since it is designed to handle highly dynamic targets and provides competitive performance under challenging motion patterns.

\paragraph{Global motion estimation and camera motion compensation.}
When the background is non-stationary (e.g., handheld camera, wind-induced scene motion), tracking-by-detection systems often benefit from estimating and compensating global frame-to-frame motion.
Frequency-domain \emph{phase correlation} estimates global translation efficiently and is widely used for fast image registration~\cite{kuglin1975phase}.
For more general parametric alignment, Enhanced Correlation Coefficient (ECC) maximization provides an iterative optimization objective that is robust to certain photometric changes and supports affine/Euclidean warps~\cite{evangelidis2008ecc}.
These alignment techniques are commonly used as a stabilization step before foreground differencing or association.

\paragraph{Optical flow and the iterative Lucas--Kanade family.}
The Lucas--Kanade (LK) method estimates motion by iteratively minimizing a photometric error under small displacement assumptions~\cite{lucas1981iterative}.
Subsequent analyses unify LK variants (forward-additive, inverse-compositional, different warp parameterizations) and clarify optimization choices and convergence behaviors~\cite{baker2004lucaskanade}.
In practice, LK is frequently used in a pyramidal coarse-to-fine manner to handle larger displacements, and is often paired with feature point selection (e.g., Shi--Tomasi corners) for stable tracking across frames~\cite{bouguet2002pyrLK,shi1994good}.
Motivated by these classical results, we plan to integrate an \emph{iterative LK} module to further refine inter-frame alignment and improve robustness under dynamic background.

\paragraph{Tracking-by-detection and data association.}
A common paradigm is to detect objects per frame and then associate detections into trajectories.
SORT~\cite{bewley2016sort} popularized a lightweight pipeline combining a simple motion model with IoU-based association for real-time multi-object tracking.
DeepSORT~\cite{wojke2017deepsort} improves association under occlusion by adding learned appearance metrics.
ByteTrack~\cite{zhang2022bytetrack} further strengthens association by leveraging low-score detections to recover missed objects.
Our pipeline follows a similarly pragmatic design: after motion compensation and candidate generation, we perform efficient association using overlap and geometry cues, aiming to remain robust while keeping compute low.

